\section{Limitations and claims not made}
\label{sec:limitations}

\paragraph{The numbers do not compose.}
The results of Section~\ref{sec:benchmarks} come from different
experiments at different degrees with different pipelines, and they
may not be multiplied. In particular two figures that circulated
internally as ``speedups'', a $131\times$ and an $18.9\times$, are not:
the first is the ratio $S/S_\rho$ of a meet-in-the-middle arm at
$n=23$ (it is $131$ times \emph{slower} than rho) and also the
entry-count quotient of a quotiented ECC2K-130 table; the second is
the compact tier's cost relative to the generic floor at $n=61$
(Table~\ref{tab:tier}). Neither is a speedup of anything, and neither
appears in this paper as one. Likewise the $6.85\times$ of
Table~\ref{tab:n19} is a CPU-time ratio between two factor bases at
$n=19$, not a ratio to rho, and should not be confused with the
unrelated $6.85\times$ that a Python Frobenius micro-benchmark happens
to report at $n=131$.\art{experiments/sage-ic-campaign/onb-frobenius-20260924/RESULT.md}

\paragraph{No end-to-end advance.}
At every degree where a complete, verified pipeline was priced against
the matched signed-Frobenius rho the index calculus is slower:
$12$--$13\times$ in operations at $n=61$, $1.3$--$23\times$ in isolated
CPU at $n=37$--$53$, $300\times$ at the best $39$-bit rung of the
boundary ledger. The one exploratory observation in the other direction
($n=83$, median $8.2\times$ in online wall) is unisolated and is not a
claim. The paper claims no attack, no asymptotic statement and nothing
about ECC2K-130 or any deployed curve.

\paragraph{The canonical form and the zero-gap fact are prior art.}
Section~\ref{sec:related}: the least Frobenius conjugate in normal-basis
representation is the normalization of \citet{bailey2009}, and
Lemma~\ref{lem:zerogap} is a standard property of necklaces
\citep{ruskeysawada1999,gabricsawada2026}. The contributions are the
implementations of the shortcut and their exactness proofs, the
layouts and the measurements.

\paragraph{Isolation.}
Of the timing rows, only the laned-key result, the cold panel of
Table~\ref{tab:cold} and the key-kernel exploration carry an
isolation receipt. The zero-gap timings, the equal-memory fold, the
probe shapes, the $n=19$ Gr\"obner comparison and the cold-job pair
tables are exploratory; their deterministic counts (rotations,
comparisons, entries, bytes, instructions) are not.

\paragraph{Scope of the proofs.}
Lemma~\ref{lem:free} and Theorem~\ref{thm:canon} assume $n$ prime, which
every measured instance satisfies; the composite case is implemented
(least shift on ties) and tested on small composite $n$ but not
measured. The sign chooser of Lemma~\ref{lem:tracesign} is proved, not
deployed: the implementations compare $y$ with $y+x$ after the
rotation, which is equivalent and cheaper when $y$ is already in hand.

\paragraph{Naming is not portable.}
A normal element found by randomized search names orbits differently
from another; a folded table, a root index or a distinguished-point
log is readable only with the basis it was written in
(Section~\ref{sec:conventions}). We treat this as an artefact
requirement rather than a defect, but it is a real one.

\paragraph{What the shortcut does not fix.}
The zero-gap scan removes rotations, not branches. On a wide
out-of-order core the data-dependent exits of variants A and B
mispredict on most keys, and the exhaustive branch-free vector scan
wins inside the decomposition loop at $n=53$
(Section~\ref{sec:canon-impl}). Which variant is right depends on the
loop it is embedded in, and the right answer at $n=83$, where no
vector register holds a key, is open.
